\scanpage{77}

% top, lines 1–3
Le passage d'un hexagone à l'hexagone dual se traduit par le passage de $(\tau,\varepsilon,\varepsilon',\varphi)$ à $(\tau,\varepsilon',\deleted{\unsure{\varepsilon'}}\varepsilon,-\varphi\circ\sigma)$ où $\sigma$ est la symétrie $\varepsilon'\wedge\varepsilon\to\varepsilon\wedge\varepsilon'$.

% upper, line 4
\[
\mathrm{Aut}(h)\subset\mathrm{Aut}(\tau)\times\deleted{\unsure{\mathrm{Aut}\,\varepsilon}}\underbrace{(\pm1)^2}_{\mathrm{Aut}\,\varepsilon\times\mathrm{Aut}\,\varepsilon'}
\]

% upper, line 5
formé des triples $u,\deleted{\unsure{\alpha}}\beta,\deleted{\unsure{\beta}}\gamma$ tels que

% upper, line 6
\[
\deleted{\unsure{\alpha\beta}}\beta\gamma=\mathrm{sg}(\unsure{\tau|u})\qquad\text{ou}\quad\alpha\beta\,\mathrm{sg}(\tau)=1
\]

% upper, lines 7–8
i.e. c'est le sous-gpe de $\mathrm{Aut}\,\tau\times(\pm1)^2$ \unsure{image|formé} inverse de $(\pm1)^{3\prime}=\{(\alpha,\beta,\gamma)\in(\pm1)^3\mid\alpha\beta\gamma=1\}$.

% middle, large diagram
\diagram{\begin{tikzcd}[ampersand replacement=\&,column sep=1.1em,row sep=1.4em,cells={font=\scriptsize}]
|[alias=D]| {\substack{\tau=\text{diagonales}\\\deleted{\unsure{\text{les faces}}}}}\arrow[ddd,"3"'] \& \& |[alias=S]| \text{Sommet}\arrow[ll,"2"']\arrow[ddd,"2"] \&\\
\& |[alias=C]| \text{côtés}\arrow[ul,"2"']\arrow[ddd] \& \&\\
\&[-1em] |[alias=R]| \text{rep. sym.}\arrow[uul,"2"]\arrow[ddd] \& \& |[alias=O]| {\substack{\text{sommets orientés}\\\simeq\text{\unsure{arcs|arêtes|côtés}}\simeq\text{repères}}}\arrow[ul,"2"]\arrow[to=C,"2"]\arrow[ll]\arrow[ddd]\\
|[alias=E]| e \& \& |[alias=B]| {\substack{\varepsilon\\\text{(triangles inscrits)}}}\arrow[ll,"2"'] \&\\
\& |[alias=F]| {\substack{\varepsilon'=\varepsilon\wedge\omega\\\text{\unsure{(triangles circonscr.)|(triangles coinscrits)}}}}\arrow[ul,"2"] \& \&\\
\& |[alias=W]| {\substack{\omega=\omega(\tau)\\\text{orientations}}}\arrow[to=E,"2"'] \& \& {\substack{\varepsilon\times\omega\\\text{triangles}\\\text{inscrits orientés}}}\arrow[ll,"2"]\arrow[to=F,"2"]\arrow[to=B,"2"]
\end{tikzcd}}

% lower, lines 1–3
\uline{NB} Tous les carrés sont cartésiens, y compris les quatre carrés horizontaux de forme triangulaire

% below NB
\diagram{\begin{tikzcd}[ampersand replacement=\&,row sep=small,column sep=small]
{}\&\&{}\arrow[ll]\\
\&{}\arrow[ul]\&\\
\&\&{}\arrow[ul]\arrow[uu]
\end{tikzcd}}

% archive sheet number: 76
